\chapter*{Wstęp}\addcontentsline{toc}{chapter}{Wstęp}\markboth{Wstęp}{Wstęp}

Ta książka jest dla dorosłego, który czyta mniej, niż chciałby, i ma listę tytułów
\glqq przeczytam kiedyś\grqq{} dłuższą niż rok. Zakładam, że masz pracę, być może rodzinę
i nie masz dodatkowych godzin, które dałoby się oddać czytaniu bez zabrania ich czemuś innemu.
Nie zakładam, że chcesz zostać zawodowym czytelnikiem ani że lubisz wyzwania liczbowe.
Liczby pojawiają się tutaj, bo pozwalają sprawdzić, czy cel jest realny, zanim go wyznaczysz.

Premisa jest prosta. Rok ma pięćdziesiąt dwa tygodnie i trzysta sześćdziesiąt pięć dni.
Liczba książek, które przeczytasz, nie jest losem ani cechą charakteru. Jest wynikiem czterech
decyzji: ile minut dziennie poświęcasz czytaniu, jakie książki wybierasz, w jakim formacie je
czytasz oraz jaki nawyk pozwala wracać do lektury, gdy dzień się nie układa. Rozdziały od
pierwszego do dziesiątego idą przez te decyzje po kolei. Część I pyta o czas, część II o wybór
i format, część III o to, co zostaje w głowie, i o prowadzenie całego roku.

Skąd wziąć czas? Nie z porannego wstawania o piątej, choć u niektórych to działa. Przegląd dnia
zwykle odsłania fragmenty, rzadko cały blok. Rozdział trzeci pokazuje, jak sprawdzić to na
własnym kalendarzu, zamiast zgadywać.

Arytmetyka jest najprostszą częścią i warto policzyć ją samemu. Przyjmujemy typową książkę
na 80\,000 słów. To założenie, a nie zestawienie wszystkich książek: kryminał bywa krótszy,
biografia dłuższa. Przyjmujemy też tempo 250 słów na minutę. Również to jest założenie modelu,
a nie pomiar czytania konkretnego człowieka, w tym ciebie. Przy tych założeniach jedna typowa
książka zajmuje \val{typical.hours} godziny. Dwanaście książek w roku to 12 razy 80\,000 słów,
czyli 960\,000 słów. Dzielimy przez 250 słów na minutę i dostajemy 3840 minut; po podzieleniu
przez 365 dni wychodzi \val{need.12} minuty dziennie. Dwadzieścia cztery książki wymagają około
\val{need.24} minut dziennie, pięćdziesiąt dwie około \val{need.52} minut. Dodatek B pokazuje te
same zależności dla innych tempa i długości, a także jak policzyć własną tablicę na podstawie
własnych pomiarów.

Te liczby mają jedną wadę, którą trzeba nazwać na początku. Liczba książek nie jest miarą
wartości. Książka na 40\,000 słów może zmienić więcej niż cegła na 120\,000, a licznik zachęca do
wybierania krótkich rzeczy i porzucania długich. Cel \glqq24 książki\grqq{} nie mówi nic o tym,
czy cokolwiek z nich zostało z tobą. Rozdziały czwarty, siódmy i ósmy zajmują się tym, co
wybrać, jak czytać trudniejsze teksty i jak sprawić, by coś zostało w pamięci. Jeśli po roku
okaże się, że przeczytałeś mniej, ale lepiej wybrałeś, to jest uczciwy wynik.

Czego ta książka nie robi. Nie podaje listy stu lektur, które trzeba znać, bo taka lista mówi więcej
o autorze listy niż o tobie. Nie ocenia książek i nie mówi, które z nich warto czytać; to, czy
tytuł ci odpowiada, zależy od twojego pytania, czasu i nastroju. Nie uczy szybkiego czytania
w rozumieniu, w którym każda strona jest przelatywana z prędkością kilku tysięcy słów na minutę,
bo w rozdziale siódmym pokazuję, dlaczego takie obietnice trzeba czytać ostrożnie. Robi jedno:
zamienia pytanie \glqq ile mogę przeczytać\grqq{} na rachunek, który możesz zmienić, gdy zmienią się
twoje warunki, i na kilka decyzji, które od ciebie zależą.

Czego ta książka nie obiecuje. Nie obiecuje, że przeczytasz dwa razy więcej, bo tytuł jest pytaniem
o możliwość, a nie gwarancją. Nie obiecuje efektów dla zdrowia, pamięci ani kariery. Tam, gdzie
twierdzenie dotyczy świata poza tekstem, a nie jest policzone w tekście, oznaczam je pomarańczowym znaczkiem [?], i traktuj je wtedy jako hipotezę do sprawdzenia w źródle. Nie obiecuje też,
że każda rada zadziała u każdego. Ta wersja nie została jeszcze przeczytana przez czytelników,
więc jej ocena pozostaje nieustalona. Nie znajdziesz tu terminów w rodzaju \glqq w trzy miesiące
przeczytasz więcej\grqq{}, bo nie mam danych, które by taki termin uzasadniały.

Książkę można czytać po kolei albo wybiórczo. Jeśli problemem jest czas, zacznij od rozdziału
trzeciego. Jeśli problemem jest powrót do lektury po przerwie, zacznij od szóstego. Przy każdym
rozdziale są ćwiczenia z odpowiedziami w dodatku A i jedno zadanie na tydzień, które wymaga zrobienia
czegoś, a nie tylko przeczytania.
